
qui se vérifie d'%
après 1.10 sur les éléments @@M0@@. Pour le premier,
@@M1@@
Pour le deuxième :
@@M2@@

\item Les exemples géométriques de 2.2 seront traités plus loin.
\end{enumerate}

2.6. On peut généraliser l'exemple 2.5.a) de la façon suivante : Soit @@M3@@ un anneau commutatif à élément unité, tel que :
@@M4@@
On pose alors, pour tout @@M5@@, et tout @@M6@@, @@M7@@, soit :
@@M8@@
où la dernière identité est considérée dans @@M9@@. D'après la multiplicativité de l'exponentielle, ceci définit une structure de pré-@@M10@@-anneau sur @@M11@@.

Pour voir que c'est une structure de @@M12@@-anneau, il faut prouver les relations 2.4.2, et il suffit de les prouver dans @@M13@@, d'après 2.6.i). On est alors ramené par fonctorialité au cas @@M14@@, et aux éléments @@M15@@. Compte tenu de la définition des @@M16@@, les deux membres dans les
